\section{The problem and the result}\label{sec:introduction}

For a Borel probability measure \(\mu\) on \(\R^3\), let
\(\Pi(\mu,\mu,\mu)\) denote the probability measures on \((\R^3)^3\)
with all three coordinate marginals equal to \(\mu\). We use the
Coulomb cost
\begin{equation}\label{eq:cost}
 c(x_1,x_2,x_3)=\frac1{|x_1-x_2|}
                 +\frac1{|x_1-x_3|}
                 +\frac1{|x_2-x_3|},
\end{equation}
where \(c=+\infty\) if two coordinates coincide, and write
\[
 \E(\mu)=\inf_{\pi\in\Pi(\mu,\mu,\mu)}\int c\dd\pi.
\]
Here and below \(|\cdot|\) is the Euclidean norm.
A \emph{Monge plan} is a coupling of the form
\((\Id,T_2,T_3)_\#\mu\), where \(T_2,T_3:\R^3\to\R^3\) are
Borel maps and \((T_2)_\#\mu=(T_3)_\#\mu=\mu\).
For a Borel map \(T\), the pushforward is defined by
\(T_\#\mu(A)=\mu(T^{-1}(A))\).
For a map defined on a set of full measure, we use the same notation
for the pushforward of the measure restricted to that set.

The question is whether absolute continuity of \(\mu\), together
with \(\E(\mu)<\infty\), guarantees an optimal Monge plan.
We give a counterexample with additional regularity.

\begin{theorem}\label{thm:main}
There is a nonnegative density \(\rho\) on \(\R^3\) such that
\[
 \int_{\R^3}\rho(x)\dd x=1,\qquad
 \rho\in C_c^\infty(\R^3),\qquad
 \sqrt{\rho}\in C_c^\infty(\R^3),
\]
with the following properties. For \(\mu=\rho\dd x\), the value
\(\E(\mu)\) is finite and is attained by a coupling in
\(\Pi(\mu,\mu,\mu)\), but every pair of Borel maps
\(T_2,T_3:\R^3\to\R^3\) satisfying
\((T_2)_\#\mu=(T_3)_\#\mu=\mu\) obeys
\[
 \int_{\R^3}c\bigl(x,T_2(x),T_3(x)\bigr)\dd\mu(x)>\E(\mu).
\]
The inequality permits the left-hand side to be infinite.
\end{theorem}

The maps in the theorem are arbitrary Borel maps. In particular, the
conclusion concerns existence of any Monge optimizer and includes
maps with no cyclic relation or regularity.
The theorem also excludes optimal symmetrizations of admissible
Monge plans, since permutation averaging preserves their marginals
and Coulomb cost.

\paragraph{History and significance.}
Coulomb multi-marginal transport arises in the strong-interaction
limit of density functional theory. In this setting the description
of optimal configurations by measure-preserving maps is usually
called the Monge or co-motion ansatz. The co-motion formulation was
developed in the strictly correlated electron literature; see Seidl,
Gori-Giorgi and Savin~\cite{SGS2007} and the account by Friesecke,
Gerolin and Gori-Giorgi~\cite{FGG2023}.
The distinction between equality of optimal values and attainment
by maps is essential. Colombo and Di Marino~\cite[author version, Theorem~1.1]{CDM2015} proved
that the Coulomb Kantorovich minimum equals the infimum over
measure-preserving cyclic maps for identical non-atomic marginals.
This provides approximating maps without asserting that a
minimizing map exists. Section~\ref{sec:infima} gives a direct proof of
equality of the two infima for the particular marginal constructed here.

For two particles, Cotar, Friesecke and
Kl{\"u}ppelberg~\cite[arXiv version, Theorem~3.6]{CFK2013}
proved existence and uniqueness of an optimal map for absolutely
continuous marginals in every dimension. For arbitrary particle number
on the line, Colombo, De Pascale and Di Marino~\cite[Theorem~1.1]{CDDM2015}
constructed optimal cyclic maps for identical nonatomic marginals with
finite transport value. Pass~\cite{Pass2013} constructed Coulomb
optimizers that are not induced by maps for radial data; the existence
of such an optimizer does not exclude a different optimizer of Monge
form. Colombo and Stra~\cite{CS2016} gave counterexamples to the
prescribed radial quantile-cyclic constructions. Bindini, De Pascale and
Kausamo~\cite{BDK2022} subsequently constructed smooth densities for
the reduced radial cost of the planar three-particle problem for which
none of the four prescribed maps III, DDI, DID and IDD is optimal.
This radial cost is obtained by minimizing the Coulomb interaction over
angles; it is different from the Coulomb cost on the line.

The general higher-dimensional map-attainment question with at least
three particles is recorded in the 2022 preprint
of~\cite[arXiv version, Theorem~2.4 and the following discussion]{FGG2023}.
The three-particle Coulomb question in \(\R^3\) was still stated as open
in January 2026 by Friesecke~\cite[Section~3.2]{Friesecke2026}.
Theorem~\ref{thm:main} excludes every preserving Borel pair for a smooth
identical marginal. With the particle-density normalization \(n=3\rho\),
one has \(\int n=3\) and \(\sqrt n\in H^1(\R^3)\). Thus the example
belongs to the physical density class \(D_3\)
of~\cite[arXiv version, Eq.~(12)]{FGG2023}. This is a density-regularity
statement; no external potential or ground-state realization is asserted.

Our construction adapts the shared-component splitting mechanism
of Gerolin, Kausamo and Rajala~\cite{GKR2019} for the repulsive
harmonic cost \(-\sum_{i<j}|x_i-x_j|^2\).
For the Coulomb interaction, the additional task is to build two
families of triples with a common supporting potential. We obtain
these families and their supporting inequalities by local calculus.

\paragraph{Idea of the proof.}
The construction has two parts: a geometric condition that forces every
optimal plan onto two prescribed families of triples, and a mass imbalance that
prevents a preserving map from choosing between them.
Let \(B\) be a small ball about zero, and let \(e_1,e_2\) be the first
two coordinate unit vectors. We construct smooth diffeomorphisms \(Y_k,W_k\), for \(k=1,2\), from \(B\) onto small
components near \(e_k,-e_k\), respectively. The two families are
\[
 (x,Y_k(x),W_k(x)),\qquad x\in B,\quad k\in\{1,2\},
\]
together with their coordinate permutations; see
Figure~\ref{fig:components}.

The triples \((0,e_k,-e_k)\) have the same Coulomb cost and the same
central gradient. This allows the two local families to share a
supporting potential on \(B\). Completing this potential on the four
outer components gives an inequality
\(c(x_1,x_2,x_3)\ge u(x_1)+u(x_2)+u(x_3)\), with equality precisely
on the two families. Strict cost gaps exclude the other triples of distinct centers;
the Coulomb singularity excludes repeated components after the
neighborhoods have been made small enough.

For a smooth probability measure \(\nu\) on \(B\), we form the
marginal \(\mu\) by giving the central component measure \(\nu/3\)
and each outer component measure \((Y_k)_\#\nu/6\) or
\((W_k)_\#\nu/6\). Randomly choosing a branch
and permuting its three coordinates then gives an optimal coupling
with this common marginal. Every optimizer must saturate the supporting
inequality. In particular, on the central component, a deterministic
map must choose one of the four branch images. If it chooses a branch
\(H\) on a Borel set \(S\subset B\), it sends mass \(\nu(S)/3\)
into \(H(S)\), which has marginal mass only \(\nu(S)/6\).
This contradicts preservation on every selection set of positive mass.
The exact image-measure relation, rather than the component totals
alone, excludes arbitrarily irregular choices of branch.

\begin{figure}[!b]
\centering
\begin{tikzpicture}[>=Stealth, every node/.style={font=\small},
 component/.style={draw,rounded corners,minimum width=2.25cm,minimum height=.85cm,align=center}]
\node[component,fill=black!4] (B) at (0,0) {$B$ near $0$\\mass $1/3$};
\node[component,draw=blue!60!black,fill=blue!4] (Y1) at (3.8,0)
 {$Y_1(B)$ near $e_1$\\mass $1/6$};
\node[component,draw=blue!60!black,fill=blue!4] (W1) at (-3.8,0)
 {$W_1(B)$ near $-e_1$\\mass $1/6$};
\node[component,draw=orange!60!black,fill=orange!5] (Y2) at (0,1.9)
 {$Y_2(B)$ near $e_2$\\mass $1/6$};
\node[component,draw=orange!60!black,fill=orange!5] (W2) at (0,-1.9)
 {$W_2(B)$ near $-e_2$\\mass $1/6$};
\draw[->,blue!60!black] (B) -- node[above] {$Y_1$} (Y1);
\draw[->,blue!60!black] (B) -- node[above] {$W_1$} (W1);
\draw[->,orange!60!black] (B) -- node[right] {$Y_2$} (Y2);
\draw[->,orange!60!black] (B) -- node[right] {$W_2$} (W2);
\end{tikzpicture}
\caption{Schematic positions of the five disjoint components, with their
marginal masses. Every optimal plan is concentrated on triples joining
a central point to its two images with the same subscript. The arrows
denote local diffeomorphisms, not straight transport trajectories.
For each indicated map \(H\) and Borel set \(S\subset B\),
\(\mu(H(S))=\mu(S)/2\).}
\label{fig:components}
\end{figure}


Sections~\ref{sec:splitting}--\ref{sec:proof} establish the mass obstruction,
construct and certify the two local families, and form the smooth marginal
and an explicit optimizer. Section~\ref{sec:infima} constructs preserving
maps whose costs approach that optimum.
Section~\ref{sec:riesz} extends the construction to every inverse-power
Riesz cost in every Euclidean dimension at least two.
